% GENERATED FILE. Edit canonical lessons, then run npm run generate:books.
% canonical-source-hash: fnv1a64:470da1722f12e28c
% canonical-lessons: PA-C252-potli, PA-C252-pand, PA-C252-van, PA-C252-takhta, PA-C252-parda, PA-R252-review-246-252, PA-R252-review-final-252

\chapter{Bundles and Curtains}
\label{ch:pa-a2-bundles-and-curtains}
\hlchaptermodality{252}

\begin{chapteropening}
\textbf{By the end of this chapter:} \emph{I can name more things around the house.}

Complete the last lesson of chapter 252: I can name more things around the house.
\end{chapteropening}

\section[\texorpdfstring{\pa{ਪੋਟਲੀ} (poṭlī)}{ਪੋਟਲੀ (poṭlī)}]{\pa{ਪੋਟਲੀ} (poṭlī) — a small bundle}

\label{lesson:PA-C252-potli}



\begin{quote}
Before the new one: say the Punjabi for a bread basket, then the Punjabi for copper.
\end{quote}

\subsection*{You'll want to know: \pa{ਪੋਟਲੀ}}
\textbf{\pa{ਪੋਟਲੀ}} — \emph{poṭlī} — \textquotedblleft{}a small bundle\textquotedblright{}.

\begin{quote}
\pa{ਦਾਦੀ} \pa{ਦੀ} \pa{ਪੋਟਲੀ} \pa{ਵਿੱਚ} \pa{ਗੁੜ} \pa{ਹੈ।}
\end{quote}

\begin{quote}
\emph{dādī dī poṭlī vich guṛ hai.}
\end{quote}

\begin{quote}
There is jaggery in my grandmother's little bundle.
\end{quote}

A feminine noun.

\begin{cousinweb}
\textbf{\pa{ਪੋਟਲੀ}} (\emph{poṭlī}) is inherited from Sanskrit \emph{pōṭṭalī}, \textquotedblleft{}bundle\textquotedblright{}, a late word that CDIAL sets under a reconstructed \emph{*pōṭṭa}, \textquotedblleft{}bundle\textquotedblright{}; CDIAL lists Punjabi \emph{poṭlī}, \textquotedblleft{}small bag\textquotedblright{} (CDIAL 8396).
\end{cousinweb}

\begin{grammarlens}[title={What you've built}]
One more word for this chapter.
\end{grammarlens}

\subsection*{Your turn}
\begin{itemize}
\raggedright
  \item \emph{Say it:} \emph{poṭlī}
  \item \emph{Say it:} \emph{poṭlī}, once more
  \item \emph{Say it:} \emph{dādī dī poṭlī vich guṛ hai.} — There is jaggery in my grandmother's little bundle.
  \item \emph{Recall:} say the Punjabi for a bread basket, then the Punjabi for copper, then say \emph{poṭlī} again
\end{itemize}

\begin{tcolorbox}[breakable,colback=teal!4,colframe=teal!35!black,title={Before you move on}]
What does \pa{ਪੋਟਲੀ} mean? (\textquotedblleft{}A small bundle\textquotedblright{}.) Say it once more.
\end{tcolorbox}
\section[\texorpdfstring{\pa{ਪੰਡ} (paṇḍ)}{ਪੰਡ (paṇḍ)}]{\pa{ਪੰਡ} (paṇḍ) — a large bundle (carried on the head)}

\label{lesson:PA-C252-pand}



\begin{quote}
Before the new one: say the Punjabi for copper, then the Punjabi for a small bundle.
\end{quote}

\subsection*{You'll want to know: \pa{ਪੰਡ}}
\textbf{\pa{ਪੰਡ}} — \emph{paṇḍ} — \textquotedblleft{}a large bundle (carried on the head)\textquotedblright{}.

\begin{quote}
\pa{ਉਸ} \pa{ਦੇ} \pa{ਸਿਰ} \pa{ਉੱਤੇ} \pa{ਕੱਪੜਿਆਂ} \pa{ਦੀ} \pa{ਪੰਡ} \pa{ਹੈ।}
\end{quote}

\begin{quote}
\emph{us de sir utte kappṛiāṁ dī paṇḍ hai.}
\end{quote}

\begin{quote}
There is a bundle of clothes on her head.
\end{quote}

A feminine noun.

\begin{cousinweb}
\textbf{\pa{ਪੰਡ}} (\emph{paṇḍ}) is inherited from Sanskrit \emph{piṇḍa}, \textquotedblleft{}lump, clod, piece\textquotedblright{}; CDIAL lists Punjabi \emph{paṇḍ}, \textquotedblleft{}bundle\textquotedblright{} (CDIAL 8168).
\end{cousinweb}

\begin{grammarlens}[title={What you've built}]
One more word for this chapter.
\end{grammarlens}

\subsection*{Your turn}
\begin{itemize}
\raggedright
  \item \emph{Say it:} \emph{paṇḍ}
  \item \emph{Say it:} \emph{paṇḍ}, once more
  \item \emph{Say it:} \emph{us de sir utte kappṛiāṁ dī paṇḍ hai.} — There is a bundle of clothes on her head.
  \item \emph{Recall:} say the Punjabi for copper, then the Punjabi for a small bundle, then say \emph{paṇḍ} again
  \item \emph{Write it:} copy \textbf{\pa{ਪੰਡ}} once, checking the whole word
\end{itemize}

\begin{tcolorbox}[breakable,colback=teal!4,colframe=teal!35!black,title={Before you move on}]
What does \pa{ਪੰਡ} mean? (\textquotedblleft{}A large bundle (carried on the head)\textquotedblright{}.) Say it once more.
\end{tcolorbox}
\section[\texorpdfstring{\pa{ਵਾਣ} (vāṇ)}{ਵਾਣ (vāṇ)}]{\pa{ਵਾਣ} (vāṇ) — cot rope; coarse rope for weaving a cot}

\label{lesson:PA-C252-van}



\begin{quote}
Before the new one: say the Punjabi for a small bundle, then the Punjabi for a large bundle.
\end{quote}

\subsection*{You'll want to know: \pa{ਵਾਣ}}
\textbf{\pa{ਵਾਣ}} — \emph{vāṇ} — \textquotedblleft{}cot rope; coarse rope for weaving a cot\textquotedblright{}.

\begin{quote}
\pa{ਮੰਜੇ} \pa{ਦਾ} \pa{ਵਾਣ} \pa{ਪੁਰਾਣਾ} \pa{ਹੈ।}
\end{quote}

\begin{quote}
\emph{manje dā vāṇ purāṇā hai.}
\end{quote}

\begin{quote}
The rope of the cot is old.
\end{quote}

A masculine noun.

\begin{cousinweb}
\textbf{\pa{ਵਾਣ}} (\emph{vāṇ}) is inherited from a reconstructed \emph{*vāṇa}, \textquotedblleft{}a sort of rush, twisted grass rope\textquotedblright{}; CDIAL lists Punjabi \emph{vāṇ}, \textquotedblleft{}rush for rope-making, coarse twine\textquotedblright{} (CDIAL 11483).
\end{cousinweb}

\begin{grammarlens}[title={What you've built}]
One more word for this chapter.
\end{grammarlens}

\subsection*{Your turn}
\begin{itemize}
\raggedright
  \item \emph{Say it:} \emph{vāṇ}
  \item \emph{Say it:} \emph{vāṇ}, once more
  \item \emph{Say it:} \emph{manje dā vāṇ purāṇā hai.} — The rope of the cot is old.
  \item \emph{Recall:} say the Punjabi for a small bundle, then the Punjabi for a large bundle, then say \emph{vāṇ} again
  \item \emph{Write it:} copy \textbf{\pa{ਵਾਣ}} once, checking the whole word
\end{itemize}

\begin{tcolorbox}[breakable,colback=teal!4,colframe=teal!35!black,title={Before you move on}]
What does \pa{ਵਾਣ} mean? (\textquotedblleft{}Cot rope; coarse rope for weaving a cot\textquotedblright{}.) Say it once more.
\end{tcolorbox}
\section[\texorpdfstring{\pa{ਤਖ਼ਤਾ} (takhtā)}{ਤਖ਼ਤਾ (takhtā)}]{\pa{ਤਖ਼ਤਾ} (takhtā) — a wooden board}

\label{lesson:PA-C252-takhta}



\begin{quote}
Before the new one: say the Punjabi for a large bundle, then the Punjabi for cot rope.
\end{quote}

\subsection*{You'll want to know: \pa{ਤਖ਼ਤਾ}}
\textbf{\pa{ਤਖ਼ਤਾ}} — \emph{takhtā} — \textquotedblleft{}a wooden board\textquotedblright{}.

\begin{quote}
\pa{ਮੈਂ} \pa{ਤਖ਼ਤੇ} \pa{ਉੱਤੇ} \pa{ਭਾਂਡੇ} \pa{ਰੱਖਦਾ} \pa{ਹਾਂ।}
\end{quote}

\begin{quote}
\emph{maiṁ takhte utte bhāṇḍe rakkhdā hāṁ.}
\end{quote}

\begin{quote}
I keep the pots on a wooden board.
\end{quote}

A masculine noun.

\begin{cousinweb}
\textbf{\pa{ਤਖ਼ਤਾ}} (\emph{takhtā}) is borrowed: it is usually traced to Persian \emph{takhta}, \textquotedblleft{}a board, a plank\textquotedblright{}.
\end{cousinweb}

\begin{grammarlens}[title={What you've built}]
One more word for this chapter.
\end{grammarlens}

\subsection*{Your turn}
\begin{itemize}
\raggedright
  \item \emph{Say it:} \emph{takhtā}
  \item \emph{Say it:} \emph{takhtā}, once more
  \item \emph{Say it:} \emph{maiṁ takhte utte bhāṇḍe rakkhdā hāṁ.} — I keep the pots on a wooden board.
  \item \emph{Recall:} say the Punjabi for a large bundle, then the Punjabi for cot rope, then say \emph{takhtā} again
  \item \emph{Write it:} copy \textbf{\pa{ਤਖ਼ਤਾ}} once, checking the whole word
\end{itemize}

\begin{tcolorbox}[breakable,colback=teal!4,colframe=teal!35!black,title={Before you move on}]
What does \pa{ਤਖ਼ਤਾ} mean? (\textquotedblleft{}A wooden board\textquotedblright{}.) Say it once more.
\end{tcolorbox}
\section[\texorpdfstring{\pa{ਪਰਦਾ} (pardā)}{ਪਰਦਾ (pardā)}]{\pa{ਪਰਦਾ} (pardā) — a curtain}

\label{lesson:PA-C252-parda}



\begin{quote}
Before the new one: say the Punjabi for cot rope, then the Punjabi for a wooden board.
\end{quote}

\subsection*{You'll want to know: \pa{ਪਰਦਾ}}
\textbf{\pa{ਪਰਦਾ}} — \emph{pardā} — \textquotedblleft{}a curtain\textquotedblright{}.

\begin{quote}
\pa{ਖਿੜਕੀ} \pa{ਉੱਤੇ} \pa{ਇੱਕ} \pa{ਨੀਲਾ} \pa{ਪਰਦਾ} \pa{ਹੈ।}
\end{quote}

\begin{quote}
\emph{khiṛkī utte ikk nīlā pardā hai.}
\end{quote}

\begin{quote}
There is a blue curtain at the window.
\end{quote}

A masculine noun.

\begin{cousinweb}
\textbf{\pa{ਪਰਦਾ}} (\emph{pardā}) is borrowed: it is usually traced to Persian \emph{parda}, \textquotedblleft{}a curtain, a veil\textquotedblright{}.
\end{cousinweb}

\begin{grammarlens}[title={What you've built}]
One more word for this chapter.
\end{grammarlens}

\subsection*{Your turn}
\begin{itemize}
\raggedright
  \item \emph{Say it:} \emph{pardā}
  \item \emph{Say it:} \emph{pardā}, once more
  \item \emph{Say it:} \emph{khiṛkī utte ikk nīlā pardā hai.} — There is a blue curtain at the window.
  \item \emph{Recall:} say the Punjabi for cot rope, then the Punjabi for a wooden board, then say \emph{pardā} again
  \item \emph{Write it:} copy \textbf{\pa{ਪਰਦਾ}} once, checking the whole word
\end{itemize}

\begin{tcolorbox}[breakable,colback=teal!4,colframe=teal!35!black,title={Before you move on}]
What does \pa{ਪਰਦਾ} mean? (\textquotedblleft{}A curtain\textquotedblright{}.) Say it once more.
\end{tcolorbox}
\section[(dialogue)]{Looking back: from Beams and Roofs to Bundles and Curtains}

\label{lesson:PA-R252-review-246-252}



\begin{quote}
Say the Punjabi for wood and a rafter.
\end{quote}

\subsection*{Your turn}
Say these aloud:

\begin{itemize}
\raggedright
  \item wood, a rafter, a thatched roof, a roof gutter, a high shelf
  \item a wall niche, a small window, a corner of a room, a low stool, a trunk
  \item a covered entrance, a door frame, a door bolt, a chain, a wall peg
  \item to light a fire, to blow on a fire, smoke, charcoal, an ember
  \item a clay cooking pot, a large pan, a metal jug, a water pitcher, brass
  \item a wooden spoon, a hanging net, a flat basket, a bread basket, copper
  \item a small bundle, a large bundle, cot rope, a wooden board, a curtain
\end{itemize}

\begin{tcolorbox}[breakable,colback=teal!4,colframe=teal!35!black,title={Before you move on}]
Wood? (\textbf{\pa{ਲੱਕੜ}}, \emph{lakkaṛ}.) A rafter? (\textbf{\pa{ਕੜੀ}}, \emph{kaṛī}.) A thatched roof? (\textbf{\pa{ਛੱਪਰ}}, \emph{chhappar}.) A roof gutter? (\textbf{\pa{ਪਰਨਾਲਾ}}, \emph{parnālā}.) A high shelf? (\textbf{\pa{ਪਰਛੱਤੀ}}, \emph{parchhattī}.) A wall niche? (\textbf{\pa{ਆਲਾ}}, \emph{ālā}.) A small window? (\textbf{\pa{ਬਾਰੀ}}, \emph{bārī}.) A corner of a room? (\textbf{\pa{ਖੂੰਜਾ}}, \emph{khūnjā}.) A low stool? (\textbf{\pa{ਪੀੜ੍ਹਾ}}, \emph{pīṛhā}.) A trunk? (\textbf{\pa{ਸੰਦੂਕ}}, \emph{sandūk}.) A covered entrance? (\textbf{\pa{ਡਿਉਢੀ}}, \emph{ḍiuḍhī}.) A door frame? (\textbf{\pa{ਚੌਖਟ}}, \emph{chaukhaṭ}.) A door bolt? (\textbf{\pa{ਅਰਲ}}, \emph{aral}.) A chain? (\textbf{\pa{ਸੰਗਲ}}, \emph{saṅgal}.) A wall peg? (\textbf{\pa{ਖੂੰਟੀ}}, \emph{khūṇṭī}.) To light a fire? (\textbf{\pa{ਬਾਲਣਾ}}, \emph{bālṇā}.) To blow on a fire? (\textbf{\pa{ਫੂਕਣਾ}}, \emph{phūkṇā}.) Smoke? (\textbf{\pa{ਧੂੰਆਂ}}, \emph{dhūṁāṁ}.) Charcoal? (\textbf{\pa{ਕੋਲਾ}}, \emph{kolā}.) An ember? (\textbf{\pa{ਅੰਗਿਆਰ}}, \emph{aṅgiār}.) A clay cooking pot? (\textbf{\pa{ਹਾਂਡੀ}}, \emph{hāṇḍī}.) A large pan? (\textbf{\pa{ਕੜਾਹਾ}}, \emph{kaṛāhā}.) A metal jug? (\textbf{\pa{ਗੜਵਾ}}, \emph{gaṛvā}.) A water pitcher? (\textbf{\pa{ਗਾਗਰ}}, \emph{gāgar}.) Brass? (\textbf{\pa{ਪਿੱਤਲ}}, \emph{pittal}.) A wooden spoon? (\textbf{\pa{ਡੋਈ}}, \emph{ḍoī}.) A hanging net? (\textbf{\pa{ਛਿੱਕਾ}}, \emph{chhikkā}.) A flat basket? (\textbf{\pa{ਛਾਬਾ}}, \emph{chhābā}.) A bread basket? (\textbf{\pa{ਚੰਗੇਰ}}, \emph{chaṅger}.) Copper? (\textbf{\pa{ਤਾਂਬਾ}}, \emph{tāmbā}.) A small bundle? (\textbf{\pa{ਪੋਟਲੀ}}, \emph{poṭlī}.) A large bundle? (\textbf{\pa{ਪੰਡ}}, \emph{paṇḍ}.) Cot rope? (\textbf{\pa{ਵਾਣ}}, \emph{vāṇ}.) A wooden board? (\textbf{\pa{ਤਖ਼ਤਾ}}, \emph{takhtā}.) A curtain? (\textbf{\pa{ਪਰਦਾ}}, \emph{pardā}.)
\end{tcolorbox}
\section[(dialogue)]{Bundles and Curtains — once more}

\label{lesson:PA-R252-review-final-252}



\begin{quote}
Say the Punjabi for a small bundle and a large bundle.
\end{quote}

\subsection*{Your turn}
\begin{itemize}
\raggedright
  \item \emph{Say it:} a small bundle, a large bundle, cot rope, a wooden board, a curtain
\end{itemize}

\begin{tcolorbox}[breakable,colback=teal!4,colframe=teal!35!black,title={Before you move on}]
A small bundle? (\textbf{\pa{ਪੋਟਲੀ}}, \emph{poṭlī}.) A large bundle? (\textbf{\pa{ਪੰਡ}}, \emph{paṇḍ}.) Cot rope? (\textbf{\pa{ਵਾਣ}}, \emph{vāṇ}.) A wooden board? (\textbf{\pa{ਤਖ਼ਤਾ}}, \emph{takhtā}.) A curtain? (\textbf{\pa{ਪਰਦਾ}}, \emph{pardā}.)
\end{tcolorbox}
